% !TeX root = ../动手学强化学习笔记.tex
% 由主文档通过 \input 载入；请编译主文档。
\setcounter{chapter}{23} % 保留原书章号。
\chapter{实验方法、综合项目与论文阅读}
\label{ch:rl-experiments}

实验把模型定义、算法更新与可观察结果连接起来。一个回报曲线只能说明指定任务、预算和实现下发生了什么；可靠解释还需要正确性检查、对照条件和统计重复。本章是各阶段共同使用的方法与项目指南，环境验证、独立评价和统计报告应从最早的表格实验开始使用，梯度与网络训练检查则随相应算法引入。项目一、二对应表格方法基础路径，项目三对应函数近似与深度方法主线，项目四供相关进阶专题选读时使用；其他专题也可用各章习题中的小规模实验检验学习成果。

\section{环境规范与最小正确性验证}

训练前固定状态与观测含义、动作范围、奖励时序、初始分布、真实终止和人为截断。用短确定性轨迹检查奖励累计和掩码，再检查随机环境的采样频率与转移表。小规模任务应保留动态规划或直接线性求解的基准，使“学到了多少”可以与已知答案比较。

算法检查优先选择能手算的输入：单步 TD 的旧值与新值，DQN 的动作索引与目标，策略梯度的有限差分，GAE 的回合边界。程序能运行、损失有限和回报偶尔上升都是必要性很弱的信号，不能替代这些针对性验证。不要同时改变环境、网络和优化器后再猜测改进来自哪里。

\section{公平比较与统计报告}

训练回报来自带探索和更新的行为，独立评价回报来自冻结参数的指定策略，二者应分别报告。比较算法时明确环境交互次数、梯度更新次数、模型规划次数、运行时间及超参数搜索预算。使用不同成本的算法，只固定其中一个尺度未必是完整公平比较。

设独立训练运行数为 $N\geq2$，第 $i$ 次训练后用 $M\geq1$ 个评价回合得到均值 $\widehat J_i$。跨训练运行的总体均值与标准误估计为
\[
\bar J=\frac1N\sum_{i=1}^N\widehat J_i,\qquad
\widehat{\rm SE}=\frac{s_J}{\sqrt N},\qquad
s_J^2=\frac1{N-1}\sum_i(\widehat J_i-\bar J)^2.
\]
独立训练运行是主要重复单位；把同一模型的许多评价回合全当作独立训练会低估训练过程的不确定性。正态近似或自助区间需要说明方法和假设，样本很少或重尾明显时应谨慎解释，并报告原始运行结果。

学习曲线平滑应固定规则且保留未平滑统计，不能只挑最佳种子。若比较多个任务，要说明分数归一化、任务权重与聚合方法；不同单位的原始奖励不宜不加解释地直接平均。调参与最终报告应使用分离的评价预算，防止重复筛选造成乐观偏差。

\section{消融实验与错误诊断}

消融从一个明确机制假设出发。例如“目标网络降低移动目标造成的不稳定”应比较有无目标网络，并监测价值尺度、目标变化和评价回报，而不仅是最终最高分。其余重要因素尽量固定；若改动也改变计算量，需要补充按相同成本比较。

错误诊断按证据层次推进：首先核对环境和数据，再核对更新公式与梯度路径，然后检查覆盖、奖励尺度和优化配置。简单基线表现更好时，应先判断复杂方法是否使用了公平预算以及实现是否正确；若核验后仍落后，报告适用边界，不以预期优势替代结果。消融失败和负面结果同样提供机制信息。

\section{综合项目}

\textbf{项目一：赌博机实验台。}实现独立随机奖励、平稳与均值切换两类任务；比较样本均值、常数步长、$\varepsilon$-贪心和 UCB，输出累计奖励、伪遗憾及多次运行的离散程度。验收时能解释初始化和探索概率，并复核一条短奖励序列的更新。

\textbf{项目二：网格世界算法对照。}先公开完整模型，用动态规划得到真值；再只向学习算法开放交互接口，比较 MC、Sarsa、Q-learning 与 Dyna。分别报告价值误差、训练行为回报、最终策略回报和计算成本；至少给出一个模型变化或终止处理错误的可解释案例。

\textbf{项目三：深度算法复现。}选定一个规模适中的离散或连续任务，明确算法版本，实现 DQN、PPO 或 SAC 中的一种。先通过手算目标、梯度与掩码验证，再开展训练；围绕一个机制做消融，保存完整配置、种子、模型检查点和原始评价数据。验收依据是实现与证据可复核，不要求人为达到未经验证的分数。

\textbf{项目四：受限信息下的决策。}从局部观测、固定数据或多智能体中选一个方向，写清相比基础 MDP 增加了什么困难，构造一个原算法失效的小例子，再比较针对性改进。讨论结果究竟支持恢复信息、改善估计还是仅改变优化效率。

\section{论文阅读与报告写作}

阅读算法论文时，先确定问题与数据条件，再抄写真实优化目标和核心更新，逐个将符号对应到本书。随后辨认理论结果对模型、覆盖、可微性或逼近误差的假设，最后检查实验是否隔离了所声称的机制。名称不同不必然意味着更新不同，名称相同也可能存在关键实现版本差异。

实验报告应能回答：解决了什么问题，比较了什么方法，花费多少数据与计算，评价了哪个策略，结果有多大波动，以及在哪些条件下失效。没有运行的方案应写为实验设计，没有证明的性质应写为假设或待检验解释。由一个小型准确例子出发，逐步扩大复杂性，比从大型任务的单条曲线推断一般规律更有助于建立可靠认识。



